\textbf{COLUNI 2020 -- Questão 10}\par\medskip No período:

\par

“Eu ia alegre, no banco da frente, com o pai, gostava de estar junto dele, fosse pra o que fosse, \underline{mesmo se} nada tivesse a dizer [...].” CARRASCOZA, J.A. \textit{Aos 7 e aos 40}. Rio de Janeiro: Alfaguara, 2016. p. 90.

\par

O significado da relação subordinativa é o mesmo identificado na alternativa:\begin{enumerate}
\item “[...] da vida que \underline{se} vivia naquela cidade, com suas ruas sujas de terra vermelha e bagaço de cana [...].”
        (CARRASCOZA, J. A. \textit{Aos 7 e aos 40}. Rio de Janeiro: Alfaguara, 2016. p. 94.)
\item “O Urso era do silêncio, igual o seu Hermes, mas \underline{quando} falava, fazia a gente se melhorar, eu me lembro
       bem de suas palavras [...].” (CARRASCOZA, J. A. \textit{Aos 7 e aos 40}. Rio de Janeiro: Alfaguara, 2016. p. 52.)
\item “[...] seu espírito vagava pelo passado; \underline{tanto que}, [...] ele se lembrou do irmão e de uns amigos da
       infância, [...].” (CARRASCOZA, J. A. \textit{Aos 7 e aos 40}. Rio de Janeiro: Alfaguara, 2016. p.30.)
\item “O homem avisou o amigo, que dele se compadeceu, \underline{embora} só na superfície, Puxa, que pena!, [...].”
        (CARRASCOZA, J. A. \textit{Aos 7 e aos 40}. Rio de Janeiro: Alfaguara, 2016. p.28.)\end{enumerate}
